% Figure: drag coefficient versus freestream Mach number showing the
% transonic drag rise (drag divergence) for a conventional vs a
% supercritical airfoil.
\begin{figure}[htbp]
\centering
\begin{tikzpicture}
\begin{axis}[
  width=12cm, height=8cm,
  xlabel={freestream Mach number $M_\infty$},
  ylabel={drag coefficient $C_D$},
  xmin=0.5, xmax=1.1, ymin=0, ymax=0.09,
  axis lines=left,
  grid=major, grid style={enggray!20},
  tick label style={font=\scriptsize},
  label style={font=\small},
  legend style={font=\scriptsize, at={(0.03,0.97)}, anchor=north west},
  legend cell align=left,
]
% conventional airfoil: flat ~0.012 until ~0.72, then sharp rise
\addplot[engblue, very thick] coordinates {
  (0.5,0.012) (0.65,0.0125) (0.72,0.014) (0.76,0.020)
  (0.80,0.034) (0.85,0.055) (0.90,0.072) (0.95,0.080) (1.05,0.078)};
\addlegendentry{conventional airfoil}
% supercritical airfoil: flat to ~0.80 then rise (M_DD pushed higher)
\addplot[engred, very thick] coordinates {
  (0.5,0.013) (0.65,0.0135) (0.78,0.0145) (0.83,0.017)
  (0.87,0.028) (0.90,0.045) (0.95,0.066) (1.0,0.074) (1.05,0.073)};
\addlegendentry{supercritical airfoil}
% M_DD markers
\draw[engblue, dashed] (axis cs:0.74,0) -- (axis cs:0.74,0.018);
\node[engblue, font=\scriptsize, anchor=south] at (axis cs:0.74,0.018) {$M_{DD}\approx 0.74$};
\draw[engred, dashed] (axis cs:0.84,0) -- (axis cs:0.84,0.020);
\node[engred, font=\scriptsize, anchor=south] at (axis cs:0.86,0.021) {$M_{DD}\approx 0.84$};
\end{axis}
\end{tikzpicture}
\caption[Transonic drag divergence]{Drag coefficient versus freestream
Mach number for a conventional and a supercritical transport airfoil
at fixed lift. The drag is nearly flat through the subsonic range,
then climbs steeply once local supersonic pockets on the upper surface
terminate in shocks strong enough to thicken and separate the boundary
layer. The Mach number where the climb begins is the drag-divergence
Mach number $M_{DD}$, conventionally taken where $dC_D/dM=0.1$. The
supercritical section (red), with its flatter upper surface and aft
camber, delays and weakens the shock and pushes $M_{DD}$ from about
$0.74$ to about $0.84$, the improvement that let transport aircraft
cruise efficiently near $M=0.80$ to $0.85$. The values are
representative of the published wind-tunnel families rather than a
specific section.}
\label{fig:vol03ch13:drag-divergence}
\end{figure}
